% Part: incompleteness
% Chapter: incompleteness-provability
% Section: godels-paper

\documentclass[../../../include/open-logic-section]{subfiles}

\begin{document}

\olfileid{inc}{inp}{gop}

\olsection{ગોડેલના મૂળ લેખ સાથે સરખામણી}

ગોડેલના ૧૯૩૧ના લેખને
થોડો સમય આપવો ઉપયોગી છે.
તેની પ્રસ્તાવનામાં હમણાં
ચર્ચેલા વિચારોનું રૂપરેખાંકન
છે. લેખ પર ઉપરછલ્લી
નજર નાખીએ તોય દરેક
તબક્કે શું થાય છે તે
સરળતાથી દેખાય છે:
પહેલાં ગોડેલ ઔપચારિક
પદ્ધતિ~$P$ વર્ણવે છે
(વાક્યરચના, સ્વયંસિદ્ધિઓ,
સાબિતીના નિયમો); પછી
આદિમ પુનરાવર્તી વિધેયો
અને સંબંધો વ્યાખ્યાયિત
કરે છે; પછી $x B y$
આદિમ પુનરાવર્તી છે એવું
બતાવે છે અને આદિમ
પુનરાવર્તી વિધેયો અને
સંબંધોનું $\Th{P}$માં
નિરૂપણ થાય છે એવી
દલીલ કરે છે. પછી ઉપર
જણાવ્યા મુજબ અપૂર્ણતા
પ્રમેય સાબિત કરે છે.
વિભાગ ૩માં તેઓ બતાવે છે
કે અસાબિતિક્ષમ વિધાનને
અંકગણિતની ભાષાનું વાક્ય
લઈ શકાય છે. અહીંથી
$\beta$-લેમ્માનો ઉદ્ભવ
થાય છે. પુનરાવર્તી
વિધેયોનું $\Th{Q}$માં
નિરૂપણ થાય છે તે
બતાવતાં અનુક્રમોને
સંભાળવા આપણે પણ આ
લેમ્મા વાપર્યું હતું.
ગોડેલ આ માટે પૂરતાં હોય
એવાં ન્યૂનતમ સ્વયંસિદ્ધિઓ
અલગ તારવતા નથી; પરંતુ
હવે જાણીએ છીએ કે
$\Th{Q}$ પૂરતું છે.
છેલ્લે, વિભાગ ૪માં
તેઓ બીજા અપૂર્ણતા
પ્રમેયની સાબિતીની
રૂપરેખા આપે છે.

\end{document}
